\chapter{Kỳ vọng và những lần bị so sánh}
\markboth{Kỳ vọng và những lần bị so sánh}{Kỳ vọng và những lần bị so sánh}

\section{Sao ba mẹ hay so con với người ta?}
\begin{truyen}{Sao ba mẹ hay so con với người ta?}{Family Talk}
\chuhoa{C}{on hỏi: ``Sao ba mẹ hay so con với con nhà người ta?''}

Ba đáp: ``Vì ba muốn con nhìn thấy người giỏi để học.''

Con nói: ``Nhưng con chỉ thấy mình bị chê thôi.''

Ba lặng đi, vì ông hiểu mục đích của mình có thể tốt mà cách làm lại gây đau.

\textit{(A parent intends comparison as motivation, but the child feels only rejection. Good intention does not cancel painful impact.)}
\end{truyen}

\begin{baihoc}
So sánh dễ làm người nghe thấy mình bị thiếu, hơn là thấy mình được khích lệ.
\textit{(Comparison often makes people feel lacking more than encouraged.)}
\end{baihoc}

\begin{ghinhoanh}
\textbf{Short English to remember:}

Comparison hurts faster than it motivates.

Impact matters with intention.
\end{ghinhoanh}

\begin{tuhocnhanh}
1. Em đã từng bị so sánh ở điểm nào nhiều nhất? \textit{(Where have you been compared most often?)}

2. Có cách nào góp ý mà không phải đem người khác ra làm thước đo? \textit{(Is there a way to guide without using others as rulers?)}
\end{tuhocnhanh}
\ngancach

\section{Con đâu phải bản sao của anh chị}
\begin{truyen}{Con đâu phải bản sao của anh chị}{Family Talk}
\chuhoa{C}{on nói: ``Con đâu phải bản sao của anh chị đâu.''}

Mẹ đáp: ``Mẹ biết, nhưng mẹ sợ con thua thiệt.''

Con nói tiếp: ``Sợ con thua thiệt không có nghĩa là ép con thành người khác.''

Mẹ nhìn con, thấy câu nói đó tuy cứng nhưng rất thật.

\textit{(A child resists being measured against siblings. The parent sees that protection can slip into erasing individuality.)}
\end{truyen}

\begin{baihoc}
Yêu nhiều đứa con không có nghĩa là được phép dùng cùng một khuôn cho tất cả.
\textit{(Loving many children does not mean using the same mold for all of them.)}
\end{baihoc}

\begin{ghinhoanh}
\textbf{Short English to remember:}

Children are not copies.

Care must fit the person.
\end{ghinhoanh}

\begin{tuhocnhanh}
1. Em có đang bị ép theo khuôn nào không? \textit{(Are you being pushed into some mold?)}

2. Trong nhà, ai cần được nhìn như một cá thể riêng hơn? \textit{(Who needs to be seen more individually at home?)}
\end{tuhocnhanh}
\ngancach

\section{Khen người khác để kích con?}
\begin{truyen}{Khen người khác để kích con?}{Family Talk}
\chuhoa{M}{ẹ nói: ``Mẹ khen người ta để con có động lực thôi.''}

Con đáp: ``Nhưng nghe nhiều quá, con chỉ thấy mình vô dụng.''

Mẹ thở dài: ``Mẹ không định làm con nghĩ vậy.''

Con nói: ``Không định không có nghĩa là không xảy ra.''

\textit{(A parent uses praise of others as a tool. The child explains that the repeated message lands as personal inadequacy.)}
\end{truyen}

\begin{baihoc}
Trong gia đình, điều người nghe cảm nhận nhiều khi quan trọng không kém điều người nói định làm.
\textit{(In a family, what the listener feels matters almost as much as what the speaker intended.)}
\end{baihoc}

\begin{ghinhoanh}
\textbf{Short English to remember:}

Intent is not the whole story.

The effect matters too.
\end{ghinhoanh}

\begin{tuhocnhanh}
1. Em đang bị tổn thương bởi cách khích lệ nào? \textit{(What kind of encouragement is actually hurting you?)}

2. Trong nhà, ai cần phân biệt rõ hơn giữa mục đích và tác động? \textit{(Who needs to separate intention from impact more clearly?)}
\end{tuhocnhanh}
\ngancach

\section{Ba mẹ tự hào về con khi nào?}
\begin{truyen}{Ba mẹ tự hào về con khi nào?}{Family Talk}
\chuhoa{C}{on hỏi: ``Ba mẹ tự hào về con khi nào?''}

Ba trả lời theo thói quen: ``Khi con làm tốt.''

Con im.

Ba chợt sửa lại: ``Không, còn khi con biết nhận lỗi, biết đứng dậy, biết sống tử tế nữa.''

Con ngẩng lên, như nghe được một điều lâu nay mình còn thiếu.

\textit{(A child asks when the parents feel proud. The parent moves from performance to values, and the answer suddenly becomes more human.)}
\end{truyen}

\begin{baihoc}
Một đứa trẻ mạnh hơn khi biết rằng mình được nhìn thấy cả ở phẩm chất, không chỉ ở thành tích.
\textit{(A child grows stronger when seen not only for achievement but also for character.)}
\end{baihoc}

\begin{ghinhoanh}
\textbf{Short English to remember:}

Pride should include character.

Children need wider mirrors.
\end{ghinhoanh}

\begin{tuhocnhanh}
1. Em muốn được cha mẹ công nhận ở điểm nào ngoài điểm số? \textit{(Where do you want recognition beyond grades?)}

2. Trong nhà, niềm tự hào đang được đặt ở đâu? \textit{(Where is pride being placed in your family?)}
\end{tuhocnhanh}
\ngancach

\section{Nếu con chỉ bình thường thì sao?}
\begin{truyen}{Nếu con chỉ bình thường thì sao?}{Family Talk}
\chuhoa{C}{on hỏi nhỏ: ``Nếu con chỉ bình thường thì sao?''}

Mẹ đáp chậm: ``Thì con vẫn là con của mẹ.''

Con cười buồn: ``Nghe vậy thì dễ, nhưng nhiều lúc con không cảm thấy vậy.''

Mẹ nắm tay con: ``Vậy lỗi là ở cách mẹ thể hiện chưa đủ, không phải ở giá trị của con.''

\textit{(A child fears not being exceptional. The parent answers with belonging, then realizes belonging must be felt, not merely said.)}
\end{truyen}

\begin{baihoc}
Điều một đứa trẻ cần không phải luôn được gọi là đặc biệt. Điều nó cần hơn là không bị đối xử như vô giá trị khi chỉ bình thường.
\textit{(A child does not always need to be called special. More importantly, they must not be treated as worthless when ordinary.)}
\end{baihoc}

\begin{ghinhoanh}
\textbf{Short English to remember:}

Ordinary does not mean unworthy.

Belonging must be felt.
\end{ghinhoanh}

\begin{tuhocnhanh}
1. Em có sợ mình không đủ nổi bật không? \textit{(Are you afraid of not standing out enough?)}

2. Trong nhà, cảm giác được thuộc về có đang đủ rõ không? \textit{(Is the sense of belonging clear enough at home?)}
\end{tuhocnhanh}
\ngancach

\section{Tại sao điểm 9 vẫn bị hỏi mất 1 điểm?}
\begin{truyen}{Tại sao điểm 9 vẫn bị hỏi mất 1 điểm?}{Family Talk}
\chuhoa{C}{on đưa bài điểm 9, mẹ hỏi ngay: ``Sao mất 1 điểm?''}

Con cười chua: ``Con cũng muốn hỏi, sao được 9 mà không ai thấy vui trước đã.''

Mẹ sững lại.

Con nói tiếp: ``Lúc nào cũng nhìn phần thiếu thì con mệt lắm.''

\textit{(A strong score is met first with what is missing. The child feels unseen in effort because lack is always noticed before progress.)}
\end{truyen}

\begin{baihoc}
Muốn con tiến bộ tiếp, đôi khi phải nhìn điều con đã làm được trước khi nói về phần còn thiếu.
\textit{(To help a child improve further, sometimes we must first see what was already done well.)}
\end{baihoc}

\begin{ghinhoanh}
\textbf{Short English to remember:}

Notice effort before deficiency.

Progress also needs recognition.
\end{ghinhoanh}

\begin{tuhocnhanh}
1. Trong nhà, thành công nhỏ có được ghi nhận không? \textit{(Are small successes recognized at home?)}

2. Em có đang chỉ nhìn phần thiếu của mình không? \textit{(Are you focusing only on what is missing in yourself?)}
\end{tuhocnhanh}
\ngancach

\section{Con có phải sống hộ ước mơ của ba mẹ?}
\begin{truyen}{Con có phải sống hộ ước mơ của ba mẹ?}{Family Talk}
\chuhoa{C}{on hỏi thẳng: ``Con có phải sống hộ ước mơ của ba mẹ không?''}

Ba trả lời rất lâu mới xong: ``Ba từng nghĩ như vậy là thương con. Nhưng có lẽ đó cũng là cách ba chưa chịu buông ước mơ cũ của mình.''

Con nói nhỏ: ``Con muốn được hướng dẫn, không muốn bị sống thay.''

Ba gật đầu.

\textit{(A child asks whether they are living out a parent's unfinished dream. The parent is forced into a painful but honest self-examination.)}
\end{truyen}

\begin{baihoc}
Cha mẹ có thể trao kinh nghiệm, nhưng không nên lấy đời con để hoàn thành đời mình.
\textit{(Parents may offer experience, but should not use a child's life to complete their own.)}
\end{baihoc}

\begin{ghinhoanh}
\textbf{Short English to remember:}

Guide a child.

Do not live through them.
\end{ghinhoanh}

\begin{tuhocnhanh}
1. Em có đang mang ước mơ nào không thật sự là của mình? \textit{(Are you carrying a dream that is not truly yours?)}

2. Trong nhà, đâu là ranh giới giữa hướng dẫn và sống thay? \textit{(Where is the line between guidance and living through someone?)}
\end{tuhocnhanh}
\ngancach

\section{Ba mẹ muốn con hơn người hay hơn chính con?}
\begin{truyen}{Ba mẹ muốn con hơn người hay hơn chính con?}{Family Talk}
\chuhoa{C}{on hỏi: ``Ba mẹ muốn con hơn người ta hay hơn chính con hôm qua?''}

Mẹ nhìn con một lúc rồi nói: ``Nếu mẹ trả lời thật, có lúc mẹ muốn cả hai.''

Con cười: ``Vậy nên con mệt.''

Mẹ gật đầu: ``Chắc mẹ phải học cách chọn điều đúng hơn là điều oai hơn.''

\textit{(A child asks whether growth or superiority is the real goal. The parent realizes how easily ego slips into parenting.)}
\end{truyen}

\begin{baihoc}
Nuôi con để hơn người rất dễ làm cả nhà mệt. Nuôi con để tốt hơn chính mình bền hơn nhiều.
\textit{(Raising a child to beat others exhausts everyone. Raising them to grow beyond themselves lasts longer.)}
\end{baihoc}

\begin{ghinhoanh}
\textbf{Short English to remember:}

Growth is better than rivalry.

Ego can hide inside parenting.
\end{ghinhoanh}

\begin{tuhocnhanh}
1. Nhà mình đang theo hơn người hay theo trưởng thành? \textit{(Is your family chasing superiority or growth?)}

2. Em tự so với mình hôm qua được không? \textit{(Can you compare yourself with who you were yesterday?)}
\end{tuhocnhanh}
\ngancach

\section{Có cần con lúc nào cũng ngoan?}
\begin{truyen}{Có cần con lúc nào cũng ngoan?}{Family Talk}
\chuhoa{C}{on hỏi: ``Có cần con lúc nào cũng ngoan không?''}

Mẹ đáp: ``Mẹ muốn con biết điều, chứ không phải lúc nào cũng phải nuốt hết cảm xúc.''

Con hỏi lại: ``Vậy nếu con không đồng ý thì sao?''

Mẹ nói: ``Thì con nói khác đi, chứ đừng giấu tới lúc bùng ra.''

\textit{(A child questions the demand to always be good. The parent distinguishes respect from emotional suppression.)}
\end{truyen}

\begin{baihoc}
Một đứa trẻ ngoan quá mức đôi khi không yên, mà chỉ đang giấu rất nhiều điều.
\textit{(An overly good child is not always at peace; sometimes they are hiding too much.)}
\end{baihoc}

\begin{ghinhoanh}
\textbf{Short English to remember:}

Respect is not silent suppression.

Children need healthy disagreement too.
\end{ghinhoanh}

\begin{tuhocnhanh}
1. Em có đang phải đóng vai quá ngoan không? \textit{(Are you forced to play the role of being too good?)}

2. Trong nhà, bất đồng có được nói tử tế không? \textit{(Can disagreement be spoken respectfully at home?)}
\end{tuhocnhanh}
\ngancach

\section{Con sợ làm ba mẹ thất vọng}
\begin{truyen}{Con sợ làm ba mẹ thất vọng}{Family Talk}
\chuhoa{C}{on nói: ``Điều con sợ nhất là làm ba mẹ thất vọng.''}

Ba đáp: ``Ba cũng sợ làm con thấy ba chỉ biết thất vọng.''

Con ngước lên.

Ba nói tiếp: ``Nếu nhà mình chỉ có nỗi sợ thất vọng thì chắc sẽ chẳng ai dám sống thật nữa.''

\textit{(A child fears disappointing the parents. The parent replies with a deeper fear: becoming someone whose disappointment silences honesty.)}
\end{truyen}

\begin{baihoc}
Gia đình mạnh không phải gia đình không có thất vọng. Đó là gia đình mà thất vọng không giết mất sự chân thật.
\textit{(A strong family is not one without disappointment, but one where disappointment does not kill honesty.)}
\end{baihoc}

\begin{ghinhoanh}
\textbf{Short English to remember:}

Disappointment should not kill honesty.

Truth must stay safe at home.
\end{ghinhoanh}

\begin{tuhocnhanh}
1. Em có đang giấu điều gì vì sợ làm ai đó thất vọng? \textit{(What are you hiding because of fear of disappointment?)}

2. Nhà mình làm gì để sự thật vẫn an toàn? \textit{(What keeps truth safe in your home?)}
\end{tuhocnhanh}
\ngancach